\section{Limitations and Falsifiability}

The framework has several explicit failure conditions.

First, pure weighting does not generically improve first-person outcomes. If accessibility is independent of outcome, the covariance contribution vanishes in expectation. Second, positive correlation is weaker than the monotone conditional-accessibility condition required for first-order stochastic dominance; nonmonotone accessibility can produce CDF crossings. Third, recognition has no effect if it changes neither trajectory utility nor accessibility. Fourth, if expected accessibility is zero, the normalized first-person measure is undefined.

The classical simulations also have limited scope. They are toy models designed to test formal identities, mechanism plausibility, and failure modes. They are not evidence that Everettian branch self-location obeys the proposed accessibility rule.

The largest unresolved physical issue is the Everett-QBS bridge assumption. A physically serious version of the theory should specify how accessibility relates to decoherence, branch measure, observer persistence, and Born-rule credence. If no defensible mapping can be given, the framework remains a mathematical observer-selection and decision model rather than a physical theory of Everettian experience.

Additional open technical issues include sequential recognition, predictor misspecification, adaptive environments, accessible-measure collapse, and how policy changes interact with branch support. These are treated as research questions rather than hidden assumptions.
